\documentclass[12pt]{article}
\usepackage[utf8]{inputenc}
\usepackage[T1]{fontenc}
\usepackage[margin=1in]{geometry}
\usepackage{setspace}
\usepackage{xcolor}
\usepackage[colorlinks=true, urlcolor=blue!60!black, linkcolor=blue!60!black]{hyperref}
\setlength{\parindent}{0em}
\setlength{\parskip}{0.45em}
\renewcommand{\baselinestretch}{1.0}
\begin{document}
\begin{center}
{\Large \textit{United States v. Belmont}, 301 U.S. 324 (1937)\par}
\vspace{0.3em}
{\normalsize Case Brief --- One-Page Edition\par}
\end{center}

\textbf{Memory Jogger.} A sole executive agreement incident to recognition needs no Senate consent and overrides state policy (\href{https://supreme.justia.com/cases/federal/us/301/324/}{Justia}).

\textbf{Facts.} The Soviet Government nationalized a Russian corporation's New York deposit with banker August Belmont. On recognizing the Soviet Union in 1933, Roosevelt accepted the Litvinov Assignment of Soviet claims to the United States, without the Senate.

\textbf{Procedural History.} The lower courts dismissed the United States' suit as contrary to New York's anti-confiscation policy.

\textbf{Votes.} Reversed 9--0: Sutherland for six; Stone, with Brandeis and Cardozo, concurred in the judgment.

\textbf{Holding.} The assignment was valid without Senate consent and overrode New York law and policy.

\textbf{Reasoning.} Recognition validated the Soviet Government's acts, and the recognition power includes settling claims. Not every compact is a treaty, and in foreign affairs ``state lines disappear.'' The Constitution has no extraterritorial operation, so the Russian corporation had to look to its own government; Belmont's executors were mere stakeholders. Stone would have rested on New York law without holding that the agreement preempts it.

\textbf{Examples.} \emph{Same-Side:} a recognition-linked claims settlement. \emph{Opposite-Side:} an agreement directing state criminal courts (\textit{Medell\'in}). \emph{Fence-Sitter:} ending tort suits without a fund.

\textbf{Critique.} Preemption by sole executive agreement bypasses the Treaty Clause and Article VI, and \textit{Medell\'in} later vindicated Stone's caution.

\textbf{Key Quotation.} ``Plainly, the external powers of the United States are to be exercised without regard to state laws or policies.'' --- Justice Sutherland, 301 U.S. at 331.

\end{document}
